\documentclass[a4paper,12pt,liststotocnumbered]{scrartcl}

\usepackage[utf8x]{inputenc}
\usepackage{ngerman}
\usepackage{tabularx}
\usepackage{graphicx}
\usepackage{hyperref}
\usepackage{calc}
\usepackage{scrpage2}

\hypersetup{colorlinks=false,pdfborder=0 0 0}
\title{Bedienungsanleitung pinacotheca}
\author{Sebastian Eberlein}

\pagestyle{scrheadings}

\chead[]{}
\lohead[Webtechnologie I - T09]{Webtechnologie I - T09}
\rohead[]{}


\begin{document}

\pagestyle{empty}

\begin{figure}[h]
	\begin{center}
		\includegraphics[width=\textwidth/2]{logo}
	\end{center}
\end{figure}

\begin{tabularx}{\textwidth}{lX}
	Fach:&Webtechnologie I\\
	Semester:&Sommersemester 2007\\
	Prüfungsleistung:&Praktischer Leistungsnachweis\\
	&\\
	Teammitglieder:&Sebastian Eberlein\\
	&Sven Salzwedel\\
	Teamnummer:&T09\\
\end{tabularx}

\begin{center}
\end{center}

\begin{center}
	\Large{\textbf{Thema:} Webbasiertes Fotoalbum}\\
\end{center}
\begin{center}
	\Large{\textbf{Bedienungsanleitung}}\\
\end{center}

\begin{center}
	Datum: \today\\
\end{center}

\newpage

\pagestyle{scrheadings}

\tableofcontents

\newpage

\section{Benutzung der Anwendung}

\subsection{Startseite der Anwendung aufrufen}

Um die Startseite der Anwendung aufzurufen, muss in einem Browsers die
URL des Rechners auf der die Anwendung läuft und der konfigurierte Port
angegeben werden (http://$<$zieladresse$>$:$<$port$>$). Als Browser
können entweder der Internet Explorer Version 6.0 von Microsoft oder der
Mozilla Firefox Version 1.5 verwendet werden.  Nach Eingabe der
angegebenen URL gelangt man auf die Startseite auf der alle vorhandenen
Alben angezeigt werden. Hier ist der Albumname, die dazugehörige
Beschreibung sowie ein Link "`show Album"' um das jeweilige Album
anzusehen, in einem farblich abgegrenzten Rahmen dargestellt.

\subsection{Fotos eines Albums ansehen}

Befindet man sich auf der Startseite, gelangt man durch klicken auf den
Link "`Show Album"' auf eine Seite "`Viewing Album"' die alle Fotos des
ausgewählten Albums als Miniaturbilder zeigt. Ist das Album durch einen
Schlüssel geschützt, wird man aufgefordert den jeweiligen Schlüssel
einzugeben. Ohne Kenntnis des vergebenen Schlüssels hat man somit keine
Recht das jeweilige Album zu sehen. Durch klicken auf ein Foto kann man
sich dieses als Detailansicht anzeigen lassen und gelangt auf die Seite
"`Viewing Photo"'. In dieser Detailansicht kann man sich, ohne in die
Miniaturansicht der Fotos des Albums zurück zu gehen, das vorhergehende
oder das nächste Foto des Albums wiederum in Detailansicht anzeigen
lassen. Hierfür existieren Links als Miniaturfotos mit der Bezeichnung
"`Previous Photo"' und "`Next Photo"'. Auf der Seite "`Viewing Photo"'
ist ein Link "`Go back to album"' vorhanden, um zur Seite "`Viewing
Album"' zurück zu kehren.

\subsection{Fotos nach Tags ansehen}

Es existier auf der Startseite ein Link "`filtered by tags"'. Dieser
Link führt auf eine Seite "`Viewing Fotos by Tags"', auf der der
Betrachter nach einen oder mehreren Tags filtern kann.  Diese Filterung
geschieht dann über die gesamte Menge der vorhandenen Fotos und Album-
übergreifend. Nach der Filterung werden die betroffenen Fotos in
Miniaturansicht angezeigt und können durch direktes anklicken in
Detailansicht betrachtet werden.  Um von der Detailansicht auf die
Filterseite zurückzukommen, existiert ein Link "`Go back to
filterview"'.  Um von der Filterseite aus zurück zur Startseite zu
kommen, existiert ein Link "`Go back to Albumlist"'.

\subsection{EXIF-Daten anzeigen}

Hat man den Schritt "`Fotos eines Albums ansehen"' durchgeführt, kann
man sich in der Detailansicht durch klicken auf den Link "`Show
additional metadata"' die EXIF-Daten\cite{exifwp} zum jeweiligen
Foto anzeigen lassen. Dieser Link existier jedoch nur, wenn das Foto
EXIF-Daten besitzt.

\subsection{Foto kommentieren}

Hat man den Schritt "`Fotos eines Albums ansehen"' durchgeführt, ist es
dem Betrachter gewährt ein Kommentar zu einem Foto hinzu zu fügen,
welcher dann von jedem anderen Betrachter sichtbar ist. Hierfür ist eine
Eingabemaske vorhanden, in dem man seinen Namen und den Kommentar
angeben und über den Klick auf den Button "`Add comment"' abschicken
kann.

\subsection{Album als Diashow anzeigen}

Auf der Seite "`Viewing Album"' existiert ein Link "`View all as
Diashow"' um die Fotos des Albums als Diashow durchlaufen zu lassen. Die
Diashow startet beim ersten Foto des Albums und endet mit dem letzten.
Außerdem wird man mit einer Meldung darauf aufmerksam gemacht, dass die
Diashow beendet ist. Bei der Diashow existier ein Link "`Go back to
album"' um während oder nach der Diashow zur Seite "`Viewing Album"'
zurück zu kehren.

\newpage

\section{Benutzung der Administrationsoberfläche}

\subsection{Administrationsoberfläche aufrufen}

Um auf die Administrationsoberfläche zu gelangen muss in einem Browsers
die URL des Rechners auf der die Anwendung läuft und der konfigurierte
Port angegeben werden (https://$<$zieladresse$>$:$<$port$>$/admin/).
Nachdem man die Adresse eingegeben hat, wird man aufgefordert sich als
Administrator zu authentifizieren. Nach dieser Benutzer- und Passwort-
abfrage gelangt man auf die Startseite der Administrationsoberfläche.

\subsection{Neues Album anlegen, bearbeiten oder löschen}

Die Startseite beinhaltet eine Eingabemaske "`Add an Album"' um ein
neues Album anzulegen und dessen Beschreibung anzugeben. Mit dem Button
"`ADD"' wird ein leeres Album mit dem angegebenen Namen und der
Beschreibung erzeugt. Das leere neue Album erscheint sofort in der
"`Album List"' und kann mit dem dazugehörigen Link "`Delete"' gelöscht
werden. Wird ein nicht leeres Album gelöscht, werden auch alle Bilder
des Albums mitgelöscht. Durch eine klick auf den Link "`Edit"' bei dem
dazugehörigen Album, gelangt man auf die Seite "`Edit Album"'. Hier kann
unter anderem im Bereich "`Edit Album"' der Albumname und die
Beschreibung geändert werden. Durch das Anklicken des Buttons "`Save
Changes"' werden die Änderungen übernommen.

\subsection{Bilder eines Albums hinzufügen, löschen oder bearbeiten}

Um Bilder eines Albums zu bearbeiten, muss der Link "`Edit"' für das
jeweilige Album geklickt werden. Dadurch gelangt man auf die Seite
"`Edit Album"'. Hier gibt es eine Eingabemaske "`Add Photo to Album"' in
der man die Beschreibung für das Foto und mit dem Button "`Durchsuchen"'
ein Foto auf dem lokalen Dateisystem angeben kann. Mit dem Button "`Add
Photo"' wird das angegeben Foto auf den Server geladen und sofort
angezeigt.  In dem Bereich "`Modify Photos"' existieren für jedes Bild
ein Link "`Delete"'. Mit einem Klick auf "`Delete"' kann ein Foto vom
Server gelöscht werden.  Um die Beschreibung eines Fotos zu ändern muss
auf die Miniaturansicht des Fotos geklickt werden. Auf der folgenden
Seite, muss in das Feld "`Photo Description"' die gewünschte
Beschreibung eingetragen und die Änderung mit dem Button "`Save
Changes"' übernommen werden.

\subsection{Tags erzeugen, ändern oder löschen}

Tags können auf der Seite "`Edit Tags"' erzeugt, umbenannt oder gelöscht
werden.  Um auf diese Seite zu gelangen, muss man auf der Startseite der
Administrationsoberfläche auf den Link "`Edit Tags"' geklickt werden.
Auf dieser Seite hat jeder Tag seinen Bereich um bearbeitet werden zu
können. Am Ende der Tagliste existiert eine Eingabemaske "`Add Tag"' um
einen neuen Tag zu erzeugen. Hierzu muss in das Textfeld der gewünschte
Tag-Name eingetragen werden.  Durch einen Klick auf den Button "`Add
Tag"' wird der Tag erzeugt und kann von diesen Zeitpunkt ab, verwendet
werden. Ein vorhandener Tag kann umbenannt werden, indem man in dem
dazugehörigem Textfeld den Namen ändert und mit dem Button "`Save
Changes"' übernommen werden. Entfernt man den Namen komplett au den
Textfeld und klickt den Button "`Save Changes"', wird der Tag gelöscht.

\subsection{Tags einem Foto zuordnen oder entfernen}

Es gibt zwei Möglichkeiten um Tags zuzuordnen. Einerseits können die
Tags auf der Seite "`Edit Album"' hinzugefügt werden. Hierzu muss bei
jedem Foto, welches zu dem Tag gehören soll, das Feld "`assign to Tag"'
markiert werden, der gewünschte Tag in der Auswahlbox ausgewählt werden
und mit dem Button "`Assign to Tag"' bestätigt werden.  Die Möglichkeit
Tags von Bildern wieder zu entfernen, besteht hier nicht. Dies muss
einzeln für jedes Foto gemacht werden, was im nächsten Absatz erklärt
wird.  Die zweite Möglichkeit um Tags zuzuordnen, besteht auf der Seite
"`Edit Photo"'. Genauer gesagt ist dies eine Kombination aus Hinzufügen
und Entfernen von Tags .  Im Bereich "`Edit Photo Data"' gibt es zwei
Boxen. Eine "`Assigned Tags"' und eine "`Not Assigned Tags"'. Mit den
Pfeilbuttons wählt man aus welche Tags man haben will. Diese müssen dann
in "`Assigned Tags"' stehen. Die Tags, die in "`Not Assigned Tags"'
stehen, werden dem Bild nicht zugeordnet oder falls sie dem Bild
zugeordnet waren, wieder entfernt.  Die Änderungen müssen mit dem Button
"`Save Changes"' übernommen werden.

\subsection{Kommentare löschen}

Auf der Seite "`Edit Photo"' können Kommentare gelöscht werden in dem
man bei dem jeweiligen Kommentar auf den Link "`delete"' klickt. Der
gewünschte Kommentar wird gelöscht und ist sowohl für den Administrator
sowie für den Betrachter nicht mehr sichtbar.

\subsection{Schlüssel für Album generieren}

Der Administrator kann festlegen ob ein Album für jeden Betrachte
gesehen werden darf, oder ob er für ein Album einen Schlüssel vergibt.
Damit kann er festlegen, dass nur derjenige das Album anschauen darf,
der den dazugehörigen Schlüssel hat. Um eine Schlüssel zu generieren,
muss man auf der Seite "`Edit Album"' im Bereich "`Manage Album Key"'
auf den Button "`Generate New Key"' klicken. Von diesem Zeitpunkt an,
muss der Betrachter zum ansehen des bestimmten Albums, genau diesen
Schlüssel angeben. Dies kann wieder rückgängig gemacht werden indem man
bei dem gewünschten Album auf "`Delete Key"' klickt.

\newpage
\addcontentsline{toc}{section}{Literatur}
\begin{thebibliography}{}
	\bibitem{exifwp} Wikipedia Beitrag zu EXIF

		\url{http://de.wikipedia.org/wiki/EXIF}
\end{thebibliography}

\end{document}
